%==============================================================================
%  Cover Letter -- Paper 2, nop SECURITY AND PRIVACY (Wiley)
%
%  Khac ban gui IJNM o hai cho, va ca hai deu co chu y:
%    1. KHONG nhac toi viec bai tung bi IJNM tu choi. Khong tap chi nao bat
%       khai chuyen do, va viet vao day la tu mo bai bang mot tin xau. Nghia
%       vu khai bao that su la KHONG NOP DONG THOI -- va bai thoa.
%    2. CO khai companion dang review o IEEE TETC. Cai nay BAT BUOC: bai
%       trich no nhu muc [2], lay do rong mach tu no, va no chua xuat ban
%       nen bien tap khong tra duoc.
%
%  Compile tren Overleaf: pdfLaTeX.
%==============================================================================
\documentclass[11pt]{letter}
\usepackage[utf8]{inputenc}
\usepackage[T1]{fontenc}
\usepackage[margin=1in]{geometry}
\usepackage{newtxtext}                 % Times, dong bo voi bai bao
\usepackage[hidelinks]{hyperref}

\signature{%
  Quan Tran Anh Vo\\
  (on behalf of all co-authors)\\
  University of Science and Technology -- The University of Danang\\
  Danang 550000, Vietnam\\
  ORCID 0009-0000-9420-1767\\
  \href{mailto:votrananhquan1111@gmail.com}{votrananhquan1111@gmail.com}}

\address{%
  Quan Tran Anh Vo\\
  University of Science and Technology\\
  The University of Danang\\
  Danang 550000, Vietnam}

\begin{document}

\begin{letter}{%
  The Editor-in-Chief\\
  \textit{Security and Privacy}}

\opening{Dear Editor-in-Chief,}

We submit for your consideration the manuscript entitled \textbf{``A Reliability
and Calibration Benchmark of QSVM Against Strong Tabular Learners on NSL-KDD and
UNSW-NB15''} as an Original Paper in \textit{Security and Privacy}.

\medskip
\noindent\textbf{What the paper does.}
An intrusion detector is deployed behind a threshold, so the quantity that governs
operations is a probability rather than a label. Quantum-kernel classifiers,
however, are almost always evaluated on accuracy alone. This paper asks the other
half of the question: when such a model reports a probability, how far can that
number be trusted?

We find first that the natural way to pose the question is unavailable. Restricting
a binned calibration estimate to the rare attack classes yields a quantity that is
identically one minus mean confidence, because every record in that subset carries
the same label. It measures confidence on known attacks and cannot detect
miscalibration at all. We verify the identity to $2.2\times10^{-16}$ across every
model and run. Several published evaluations report exactly this quantity as a
calibration error.

Measured where both classes are present, over ten paired runs with Holm-corrected
Wilcoxon tests and on two corpora at two circuit widths, a four-qubit quantum kernel
is better calibrated than both Random Forest and XGBoost in all eight comparisons,
and is not better than a small multilayer perceptron anywhere. Post-hoc
recalibration helps the quantum kernel and no other model. One measured property of
the score distributions accounts for that and for two further results: the tree
ensembles emit a distinct probability for only $2.3\%$ of test records against
$73.6\%$ for the quantum kernel, and a monotone recalibrator cannot separate records
that already share a score.

\medskip
\noindent\textbf{Why \textit{Security and Privacy}.}
The contribution is a reliability result about intrusion detectors rather than an
accuracy claim about quantum computing. It concerns whether a deployed detector's
probabilities can be acted on, which we believe is squarely within the journal's
scope. We report the negative findings as prominently as the positive ones: of 80
Holm-corrected comparisons, 25 favour the quantum kernel and 30 favour a classical
baseline, and the quantum kernel does not rank better than any baseline anywhere.

\medskip
\noindent\textbf{Reproducibility.}
Every number in the manuscript is computed from released artifacts and reaches the
text through a generated macro; no value is typed by hand. The artifacts, the
figure and table scripts, and an audit script that recomputes each reported number
and fails if any disagrees are archived at
\href{https://doi.org/10.5281/zenodo.22893683}{https://doi.org/10.5281/zenodo.22893683}.

\medskip
\noindent\textbf{Related submission.}
A companion manuscript by the same group, \textit{``NISQ-aware quantum kernel SVM
for network intrusion detection: a regime-specific benchmark on NSL-KDD''}, is
currently under review at \textit{IEEE Transactions on Emerging Topics in
Computing} (manuscript TETC-2026-05-0252). It is cited as reference [2] of the
present manuscript.

The two do not overlap. The companion studies classification accuracy and the
regime in which a quantum kernel is competitive; the present manuscript studies
calibration and reliability and makes no accuracy claim. The present work takes
only the embedding width from the companion, which we state explicitly in Section
II-A so that the two results remain separable, and re-derives everything else. We
mention this so that the relationship is on the record rather than discovered
later.

\medskip
\noindent\textbf{Declarations.}
The manuscript is original, is not under consideration elsewhere, and has not been
published. All authors have approved the submission and the author contribution
statement. The study used only public benchmark datasets and involved no human
participants and no personally identifiable information. The authors declare no
conflict of interest.

\medskip
We hope the manuscript is of interest to the readership, and we thank you for your
time.

\closing{Yours sincerely,}

\end{letter}
\end{document}
